\section{Related Work}
\label{sec:related-work}

\subsection{Agent Security and Governance}

Wang et al.~\cite{wang2024agents} survey LLM-based autonomous agents,
including their construction and evaluation. Ramchurn et
al.~\cite{ramchurn2004trust} examine trust in multi-agent interactions,
while AgentDojo~\cite{debenedetti2024agentdojo} evaluates prompt injection
attacks and defenses for tool-using agents. These works primarily focus on
agent behavior, trust, and evaluation, informing how agent interactions can
be governed. Our governance question concerns authorization decision
consistency between policy evaluation and enforcement: does granting contact
to an identified agent preserve the receiving user's intended policy?
This differs from assessing behavioral trust or resistance to prompt injection.

\subsection{Capability and Token-Based Authorization}

SPKI~\cite{ellison1999spki} represents naming and authority in certificates.
OAuth~2.0~\cite{hardt2012oauth} supports scoped service access.
Macaroons~\cite{birgisson2014macaroons} support delegation and attenuation
through contextual caveats. These designs address credential representation
and controlled delegation.

Credential correctness does not automatically establish authorization decision
correctness. Issuance must depend on a decision consistent with the intended
policy, not merely produce a credential that later passes validation. We
examine that decision-to-enforcement connection in SAGA; we propose neither
a new token format nor an attenuation mechanism.

\subsection{Formal Verification of Authentication Protocols}

ProVerif~\cite{blanchet2016proverif} and Tamarin~\cite{meier2013tamarin}
support symbolic analysis of properties including authentication correspondences
and secrecy. Their results depend on the modeled behavior, assumptions, and
queries; neither tool is intrinsically restricted to authentication.

Our distinction concerns verification scope: an authorization claim must
include authorization-relevant decision layers or justify their connection
to the model. The original queried authentication and secrecy results remain
valid in our diagnostic extension. Separately, ProVerif analyzes communication
consequences under explicit authorization-divergence assumptions. It does not
execute the Python matcher or establish implementation correctness. The
coverage question concerns the model-to-claim relationship, not tool soundness.

\subsection{Formal Verification of Authorization Policies}

Cassandra~\cite{becker2004cassandra} specifies policy evaluation and
enforcement semantics. SecPAL~\cite{becker2007secpal} connects declarative
authorization to an evaluation strategy with soundness, completeness, and
termination results under stated conditions. Margrave~\cite{fisler2005margrave}
checks policy properties and change impact for a supported subset of XACML.
These works analyze policy semantics and consistency; they also address the
relationship between policies, evaluation, and their surrounding software.

Our case study connects an observed implementation-level disagreement with a
rule-selection contract to its modeled protocol-level consequence. This is
not a competing policy language or a general evaluation-correctness result.
Concrete matcher evidence and formal consequence evidence remain distinct;
a strict semantic mapping between them is unfinished.

\subsection{Authentication--Authorization Gaps}

Authentication asks ``who are you,'' while authorization asks ``what may you
do.'' These are distinct obligations, but their connection is longstanding:
Lampson et al.~\cite{lampson1992authentication} develop authentication theory
together with distributed access control. SPKI~\cite{ellison1999spki} and
SecPAL~\cite{becker2007secpal} likewise make authority and its derivation
explicit. The distinction itself is not our novelty claim. These foundations
leave open the question of whether a particular agent implementation preserves
its intended policy semantics across protocol actions.

In SAGA, a software authorization decision can feed material release,
credential issuance, and message acceptance. We study the gap between that
decision's implementation and the actions consuming it. The required
consistency spans policy semantics, policy evaluation implementation, and
protocol enforcement. A correspondence to prior token issuance establishes
credential provenance, but does not alone establish that issuance was
permitted under the user's policy.

Our contribution relates a confirmed local authorization inconsistency to
formal analysis of its possible communication consequence under explicit
assumptions. The gate controls further distinguish the modeled release and
issuance boundaries. This connection is supported by two evidence layers,
not an executable refinement proof. We neither claim to originate the study
of authorization gaps nor report a production deployment exploit. The
remaining matcher-to-model mapping limits the cross-layer conclusion without
negating the confirmed implementation discrepancy.

% Positioning sources: paper/PAPER_STORY.md; paper/sections/01_introduction.tex;
% paper/sections/03_problem_formulation.tex; paper/sections/08_discussion.tex.
% All twelve existing references are retained; bibliography unchanged.
